\scanpage{45}
\mspage{22}
ceci est sans doute faux (cf plus bas \inserted{n\textsuperscript{o} 3}) --- il faudrait \inserted{\unsure{en}} \deleted{\unsure{donc|dans}} \inserted{renforcer} c) ci-dessus en

c') L'hom. $\mathcal E_{K'}^\circ-\mathcal E_K$ induit par $\mathcal E_{K'}-\mathcal E_K$ est anabélien (où $\mathcal E_{K'}^\circ$ est le noyau de l'hom. composé
\[
\mathcal E_{K'}^{\unsure{\circ}}\longrightarrow\Gamma_{\overline{\mathbb Q}/\mathbb Q}\xrightarrow{\substack{\chi\ \text{caractère}\\\text{cyclotomique}}}\widehat{\mathbb Z}^{*}).
\]

Mais pour voir que cette condition est \uline{nécessaire} pour que l'hom. soit géométrique, il faudrait \deleted{savoir} vérifier que \inserted{pour} tout sous-groupe ouvert $\mathcal E'$ d'un $\mathcal E_K$, le centralisateur dans $\mathcal E_K$ (non seulement de $\mathcal E'$ lui-$=$, mais même de $\mathcal E'^\circ$) est réduit à $1$ --- ce qui résulte de la démonstration du th. 1, et du fait* que pour tt sous-groupe ouvert $\Gamma'$ de $\Gamma=\Gamma_{\mathbb Q}$, le centralisateur (non seulement de $\Gamma'$, mais $=$) de $\Gamma'^\circ$ dans $\Gamma$ est réduit à $(1)$.
\margin{* à vérifier!}

Donc, la conjecture initiale ainsi \unsure{être|et} corrigée donne la

\uline{Conséquence} (conjecturale) Pour toute \deleted{\inserted{\unsure{groupe de}}} section de $\mathcal E_K$ sur \deleted{$\Gamma_{\mathbb Q}$, \unsure{soit} \illegible{1}} \inserted{un sous-groupe ouvert $\Gamma'$ de $\Gamma_{\mathbb Q}$,} de sorte que $\Gamma'$ opère (effectivement) sur $\pi_K$, on a \inserted{(\unsure{si} $K$ pas alg. sur $\mathbb Q$ i.e. $\pi_K\ne(1)$)} $\pi_K^{\Gamma'^\circ}\ne(1)$.
\margin{44}
